\par\addvspace{12pt}\Needspace{9\baselineskip}

\section{微分中值定理}

% ITEM: calculus-example-016
\solutionitem{calculus-example-016}{例5.16}
比较\(x^2\)、\(x^3\)与\(|x|\)在原点的极值和可导性。再考察\(f(x)=x\)在\([0,1]\)上的最值，说明驻点条件为什么不能替代端点比较。
\solutionanswer
\(x^2\)在原点可导、导数为零并取得严格极小值；\(x^3\)在原点可导、导数为零但没有极值；\(|x|\)在原点不可导却取得严格极小值。故内点可导极值必是驻点，但驻点未必为极值点，不可导点也可能取极值。\(f(x)=x\)在\([0,1]\)上最小值\(0\)、最大值\(1\)分别在两端取得，内部没有驻点。


% ITEM: calculus-example-017
\solutionitem{calculus-example-017}{例5.17}
对\(f(x)=x^3-x\)，在\([-1,1]\)上求Rolle定理所保证的点。再考察\(|x|\)在\([-1,1]\)上的情形，说明可导条件的作用。
\solutionanswer
两端函数值均为零，导数\(3x^2-1\)在\((-1,1)\)内的零点为\(\pm1/\sqrt3\)，都是Rolle定理保证的点。\(|x|\)虽两端值相等，却在原点不可导，其他点导数只为\(\pm1\)，没有零点；缺少内部可导条件时不能使用定理。


% ITEM: calculus-example-018
\solutionitem{calculus-example-018}{例5.18}
不求\(\sqrt{101}\)的精确值，利用\(\sqrt x\)在\([100,101]\)上的中值定理，给出\(\sqrt{101}-10\)的上下界。与在\(100\)处作一次近似比较。
\solutionanswer
中值定理给出\(\sqrt{101}-10=1/(2\sqrt\xi)\)，其中\(100<\xi<101\)，所以
\[
 \frac1{2\sqrt{101}}<\sqrt{101}-10<\frac1{20}.
\]
在\(100\)处的一次近似给出\(\sqrt{101}\approx10.05\)，中值定理还说明它偏大，并给出上述确定的上下界。


% ITEM: calculus-example-019
\solutionitem{calculus-example-019}{例5.19}
对\(x>0\)，证明
\begin{equation*}
 \frac{x}{1+x}<\ln(1+x)<x.
\end{equation*}
再设\(f'(x)=2x\)、\(f(1)=3\)，确定\(f\)。
\solutionanswer
对\(\ln t\)在\([1,1+x]\)使用中值定理，得\(\ln(1+x)=x/\xi\)，其中\(1<\xi<1+x\)，两边界随即得到。又\((f(x)-x^2)'=0\)，中值定理说明差为常数；由\(f(1)=3\)得到\(f(x)=x^2+2\)。


% ITEM: calculus-example-020
\solutionitem{calculus-example-020}{例5.20}
（2024年北京卷第20题第（2）问改编）
设\(f(x)=x+k\ln(1+x)\)，\(k\ne0\)。证明对任意\(t>0\)，曲线在\((t,f(t))\)处的切线都不经过原点。尝试用\(\ln(1+x)\)的中值定理解释这一结论。
\solutionanswer
切线在纵轴上的截距为
\[
 f(t)-tf'(t)=k\left[\ln(1+t)-\frac{t}{1+t}\right].
\]
对\(\ln x\)在\([1,1+t]\)用中值定理得\(\ln(1+t)=t/\xi>t/(1+t)\)。因\(k\ne0\)，截距非零，所以切线不经过原点。


% ITEM: calculus-example-021
\solutionitem{calculus-example-021}{例5.21}
设\(f'=f\)，证明\(e^{-x}f(x)\)为常数；若再给定\(f(0)=1\)，确定\(f\)。
\solutionanswer
\((e^{-x}f(x))'=e^{-x}(f'-f)=0\)。由中值定理，\(e^{-x}f(x)=C\)，即\(f(x)=Ce^x\)。初值给出\(C=1\)，故\(f=e^x\)。


% ITEM: calculus-exercise-07-01
\solutionitem{calculus-exercise-07-01}{习题5.7（1）}
证明\(|\arctan x-\arctan y|\le|x-y|\)。
\solutionanswer
若\(x\ne y\)，中值定理给出\(\arctan x-\arctan y=(x-y)/(1+\xi^2)\)，取绝对值即得；\(x=y\)时成立。


% ITEM: calculus-exercise-07-02
\solutionitem{calculus-exercise-07-02}{习题5.7（2）}
设\(f\)在\([0,1]\)连续、内部可导，且\(f(0)=0\)、\(1\le f'(x)\le2\)。估计\(f(1)\)及\(f(3/4)-f(1/4)\)。
\solutionanswer
分别在\([0,1]\)、\([1/4,3/4]\)用中值定理，得到\(1\le f(1)\le2\)、\(1/2\le f(3/4)-f(1/4)\le1\)。


% ITEM: calculus-exercise-07-03
\solutionitem{calculus-exercise-07-03}{习题5.7（3）}
证明\(e^x=1+x\)只有一个实根。可先用Rolle定理控制根数，再判断函数值。
\solutionanswer
\(0\)是一个根。若还有非零根\(c\)，Rolle定理会在\(0,c\)之间给出\(e^\xi-1=0\)，即\(\xi=0\)，不在两端之间，矛盾。因此根只有\(0\)。


% ITEM: calculus-exercise-07-04
\solutionitem{calculus-exercise-07-04}{习题5.7（4）}
若\(f\)在\(\mathbb R\)上可导且有周期\(T>0\)，证明任意长度为\(T\)的开区间内至少有一个\(f'\)的零点。
\solutionanswer
任取区间\((a,a+T)\)，周期性给出两端值相等，Rolle定理就在其内部给出导数零点。


% ITEM: calculus-exercise-07-05
\solutionitem{calculus-exercise-07-05}{习题5.7（5）}
设\(f\)在\(\mathbb R\)上可导。证明：\(f'\)为偶函数，当且仅当\(f(x)-f(0)\)为奇函数；\(f'\)为奇函数，当且仅当\(f\)为偶函数。结合例5.13，说明函数值\(f(0)\)的作用。
\solutionanswer
若\(f'\)偶，令\(H(x)=f(x)+f(-x)-2f(0)\)，则\(H'=f'(x)-f'(-x)=0\)，且\(H(0)=0\)，故\(f(x)-f(0)\)为奇函数。反过来对这个奇函数关系求导，即得\(f'\)偶。若\(f'\)奇，令\(K(x)=f(x)-f(-x)\)，则\(K'=f'(x)+f'(-x)=0\)，由\(K(0)=0\)得\(f\)偶；逆向仍由求导得到。
前例中\(f'\)偶只能确定减去常数后的奇性；例如\(a=b=1\)时\(f=e^{2x}-e^{-2x}-cx\)恰有\(f(0)=0\)，才直接为奇函数。


% ITEM: calculus-exercise-07-06
\solutionitem{calculus-exercise-07-06}{习题5.7（6）}
对\(0<a<b\)，证明\(\dfrac{b-a}{b}<\ln\dfrac ba<\dfrac{b-a}{a}\)，并据此估计\(\ln\dfrac{101}{100}\)。
\solutionanswer
对\(\ln x\)在\([a,b]\)用中值定理，\(\ln(b/a)=(b-a)/\xi\)，\(a<\xi<b\)，故两严格不等式成立。取\(a=100,b=101\)，得\(1/101<\ln(101/100)<1/100\)。


% ITEM: calculus-exercise-07-07
\solutionitem{calculus-exercise-07-07}{习题5.7（7）}
设\(f\)在\([a,b]\)上连续、内部可导，且\(|f'|\le K\)。若\(f(a)=f(b)=0\)，证明
\begin{equation*}
 |f(x)|\le K\min\{x-a,b-x\}\le\frac{K(b-a)}2\qquad(x\in[a,b]).
\end{equation*}
\solutionanswer
分别与两端比较，中值定理给出\(|f(x)|\le K(x-a)\)、\(|f(x)|\le K(b-x)\)，取较小者即得；两端的结论由函数值直接成立。


% ITEM: calculus-exercise-07-08
\solutionitem{calculus-exercise-07-08}{习题5.7（8）}
设区间\(I\)上\(f'\ge m>0\)，且\(f(\alpha)=0\)。证明\(|x-\alpha|\le|f(x)|/m\)（\(x\in I\)），并解释它对估计近似根误差的作用。
\solutionanswer
当\(x\ne\alpha\)时，\(f(x)=f'(\xi)(x-\alpha)\)，所以\(|f(x)|\ge m|x-\alpha|\)。它把可以算出的函数残差\(|f(x)|\)转为根的距离误差上界；\(x=\alpha\)时也成立。


% ITEM: calculus-exercise-07-09
\solutionitem{calculus-exercise-07-09}{习题5.7（9）}
（2013年北京卷理科第18题第（2）问）曲线\(y=\ln x/x\)在\((1,0)\)处的切线为\(y=x-1\)。证明除切点外，曲线都位于这条直线的下方。可先用中值定理比较\(\ln x\)与\(x-1\)。
\solutionanswer
中值定理给出\(\ln x<x-1\)（\(x\ne1\)）：\(x>1\)时由导数\(1/t<1\)，\(0<x<1\)时在\([x,1]\)比较。若\(x>1\)，\(\ln x/x<\ln x<x-1\)；若\(0<x<1\)，\(\ln x/x<(x-1)/x<x-1\)。故除\(1\)外均严格在切线下方。


% ITEM: calculus-exercise-07-10
\solutionitem{calculus-exercise-07-10}{习题5.7（10）}
设\(f\)在\([a,b]\)上连续、内部可导，且有三个不同零点。证明\(f'\)至少有两个不同零点。若\(f'\)也可导，证明它的导函数至少有一个零点。
\solutionanswer
将三个零点排成\(x_1<x_2<x_3\)，在两个相邻区间用Rolle定理，得到\(u\in(x_1,x_2),v\in(x_2,x_3)\)，满足\(f'(u)=f'(v)=0\)。若\(f'\)可导，再在\([u,v]\)用Rolle定理，得到\(f''\)的零点。


% ITEM: calculus-exercise-07-11
\solutionitem{calculus-exercise-07-11}{习题5.7（11）}
设\(f\)在\([0,1]\)上连续、内部可导，且\(f(0)=f(1)=0\)、\(f(1/2)=1\)。证明存在\(u\in(0,1/2)\)、\(v\in(1/2,1)\)，使\(f'(u)-f'(v)=4\)。
\solutionanswer
分别在\([0,1/2]\)、\([1/2,1]\)使用中值定理，得到\(f'(u)=2,f'(v)=-2\)，故差为\(4\)。


% ITEM: calculus-exercise-07-12
\solutionitem{calculus-exercise-07-12}{习题5.7（12）}
设\(f\)在\([0,1]\)上连续、内部可导，且\(f(0)=0\)、\(f(1)=1\)。证明存在\(\xi\in(0,1)\)，使\(f'(\xi)=2\xi\)。说明应先从\(f\)中减去哪个函数，再使用Rolle定理。
\solutionanswer
取\(F(x)=f(x)-x^2\)，则\(F(0)=F(1)=0\)。Rolle定理给出\(F'(\xi)=f'(\xi)-2\xi=0\)。减去的函数同时匹配两端值，并让导数成为目标中的\(2x\)。


% ITEM: calculus-exercise-07-13
\solutionitem{calculus-exercise-07-13}{习题5.7（13）}
设\(f\)将\([0,1]\)映入自身，在该区间连续、内部可导，且\(|f'|\le q<1\)。证明方程\(f(x)=x\)有且只有一个解，分别指出连续性与导数上界的作用。
\solutionanswer
令\(F(x)=f(x)-x\)。有\(F(0)\ge0,F(1)\le0\)，连续性及介值定理保证零点存在。若\(u<v\)都是不动点，中值定理会给出\(1=[f(v)-f(u)]/(v-u)=f'(\xi)\)，与\(|f'|\le q<1\)矛盾。连续性负责存在，导数上界负责唯一。


% ITEM: calculus-exercise-07-14
\solutionitem{calculus-exercise-07-14}{习题5.7（14）}
设\(f\)在\([0,1]\)上连续、内部可导，两端的函数值为零。证明存在\(\xi\in(0,1)\)，使\(f'(\xi)=f(\xi)\)。
\solutionanswer
为使导数出现\(f'-f\)，取\(F(x)=e^{-x}f(x)\)。两端值为零，Rolle定理给出\(e^{-\xi}(f'(\xi)-f(\xi))=0\)，即所求。


% ITEM: calculus-exercise-07-15
\solutionitem{calculus-exercise-07-15}{习题5.7（15）}
若区间上\(f'\ne0\)，证明\(f\)至多有一个零点。
\solutionanswer
若有两个不同零点，Rolle定理必在其间产生\(f'\)的零点，与题设矛盾。


% ITEM: calculus-example-022
\solutionitem{calculus-example-022}{例5.22 *}
一个函数在零点左侧恒为\(-1\)，右侧恒为\(1\)。无论怎样规定它在零点的值，它能否成为某个可导函数在整个区间上的导函数？与\(x^2\sin(1/x)\)的导函数比较。
\solutionanswer
不能。若它是某个可导函数的导数，Darboux定理要求它取遍\((-1,1)\)，而除零点外它只取\(\pm1\)，在零点补一个值仍会漏掉其他中间值。反之，\(x^2\sin(1/x)\)补零后的导数在零点附近振荡不连续，却确实来自可导函数，仍有导数的介值性。不连续与跳跃不是同一个性质。


% ITEM: calculus-exercise-08-01
\solutionitem{calculus-exercise-08-01}{习题5.8（1） *}
设\(f(x)=e^x-1\)（\(x<0\)），\(f(x)=ax\)（\(x\ge0\)）。分别用差商与导数极限定理，求\(f\)在零点可导的条件。
\solutionanswer
\(f\)在零点连续。差商的左极限为\(1\)，右极限为\(a\)，故可导当且仅当\(a=1\)。另一方法：零点两侧导函数分别为\(e^x\)、\(a\)，极限为\(1,a\)，用导数极限定理的单侧版本得到同样的条件。


% ITEM: calculus-exercise-08-02
\solutionitem{calculus-exercise-08-02}{习题5.8（2） *}
设区间上可导函数\(f\)的导数只取\(-1\)、\(1\)中的值。证明\(f\)必为\(x+C\)或\(-x+C\)，其中\(C\)为常数。
\solutionanswer
若导数取过两个不同值，Darboux定理要求它也取\(0\)，矛盾。因此\(f'\)恒为\(1\)或恒为\(-1\)，由中值定理得\(f=x+C\)或\(f=-x+C\)。


% ITEM: calculus-exercise-08-03
\solutionitem{calculus-exercise-08-03}{习题5.8（3） *}
若\(f\)在开区间内可导，且\(f'\)单调，证明\(f'\)连续。可以先在任意内点的两侧建立单调函数的单侧极限，再使用导数极限定理。
\solutionanswer
以内点\(c\)为例，单调性使\(f'\)在两侧有极限，且被邻近两点的有限导数值夹住，所以均有限。\(f\)在\(c\)连续，单侧导数极限定理分别说明这两个极限都等于\(f'(c)\)，故\(f'\)在\(c\)连续。


% ITEM: calculus-exercise-09-01
\solutionitem{calculus-exercise-09-01}{习题5.9 *}
对\(0<a<b\)，取\(f(x)=\ln x\)、\(g(x)=x^2\)，证明存在\(\xi\in(a,b)\)，使\(\ln(b/a)/(b^2-a^2)=1/(2\xi^2)\)，并据此给出该比值的上下界。
\solutionanswer
Cauchy中值定理给出\((\ln b-\ln a)2\xi=(b^2-a^2)/\xi\)，所以所求比值为\(1/(2\xi^2)\)。由\(a<\xi<b\)，得
\[
 \frac1{2b^2}<\frac{\ln(b/a)}{b^2-a^2}<\frac1{2a^2}.
\]


% ITEM: calculus-example-023
\solutionitem{calculus-example-023}{例5.23 *}
求\(\lim_{x\to0}(e^x-1-x)/x^2\)及\(\lim_{x\to+\infty}\ln x/x\)。再比较\(\lim_{x\to0}(\sqrt{1+x}-1)/x\)使用有理化与洛必达法则的计算。
\solutionanswer
洛必达法则连续使用两次给出第一极限\(\lim(e^x-1)/(2x)=1/2\)；第二极限由\((\ln x)'/(x)'=1/x\to0\)，得\(0\)。第三极限有理化为\(1/(\sqrt{1+x}+1)\to1/2\)；洛必达也给出\(1/[2\sqrt{1+x}]\to1/2\)。此处有理化无需额外求导，步骤更少。


% ITEM: calculus-exercise-10-01
\solutionitem{calculus-exercise-10-01}{习题5.10（1） *}
求\(\lim_{x\to0}(x-\sin x)/x^3\)及\(\lim_{x\to0}(\ln(1+x)-x)/x^2\)。
\solutionanswer
第一式连续用洛必达法则，化为\(\lim(1-\cos x)/(3x^2)=\lim\sin x/(6x)=1/6\)；第二式化为\(\lim[-x/(1+x)]/(2x)=-1/2\)。


% ITEM: calculus-exercise-10-02
\solutionitem{calculus-exercise-10-02}{习题5.10（2） *}
求\(\lim_{x\to+\infty}x e^{-x}\)与\(\lim_{x\to0^+}x\ln x\)。
\solutionanswer
第一式为\(x/e^x\)，洛必达给出\(1/e^x\to0\)。第二式写成\(\ln x/(1/x)\)，导数比为\(-x\to0\)，所以极限也为\(0\)。


% ITEM: calculus-exercise-10-03
\solutionitem{calculus-exercise-10-03}{习题5.10（3） *}
计算\(\lim_{x\to0}(1+x)^{1/x}\)。说明若用对数求导和洛必达法则，这个计算为什么只是核对已有结论，不能反过来承担本章指数、对数求导公式的最初建立。
\solutionanswer
取对数为\(\ln(1+x)/x\)，洛必达后极限为\(1\)，原极限为\(e\)。但本章对数、指数的导数已依赖于此前建立的指数极限及\(e\)的理论；此处是在这些结果基础上核对，不能把依赖它们的法则再当作它们的首次证明。


% ITEM: calculus-exercise-10-04
\solutionitem{calculus-exercise-10-04}{习题5.10（4） *}
令\(f(x)=x+\sin x\)、\(g(x)=x\)。比较\(x\to+\infty\)时原比值与导数比的极限。
\solutionanswer
原比值为\(1+\sin x/x\to1\)，导数比为\(1+\cos x\)，没有极限。这说明原比值收敛不要求导数比收敛，不能逆用洛必达法则。
\par
